% ============================================================================
% Fichier  : partie4/P04_C21_audit_securite_technique.tex
% Titre    : Auditer la Sécurité Technique de l'Infrastructure
% Auteur   : MINKA MI NGUIDJOÏ Thierry Emmanuel
% Version  : 1.0 -- novembre 2026
% ============================================================================
% Base : ACSI C16 (IN.1 à IN.5 ; matrice de flux ; quatre centres de LiZbiZ
% Douala 4/1/0, Yaoundé 4/1/0, Garoua 0/4/1, Bafoussam 0/3/2, ensemble
% 8/9/3, µ 0,40 U 0,15 H 1,46 ; flux DMZ→utilisateurs 445 et
% utilisateurs→base 1433 non autorisés). Base du lab ASI-LAB-11.
% Sources : ISO/IEC 27002:2022 (mesures 8.8, 8.9, 8.15, 8.20, 8.22, 8.24,
% numéros à vérifier dans l'édition en vigueur) ; Center for Internet
% Security, CIS Benchmarks et CIS Controls v8 ; FIRST, CVSS v3.1 ; NIST
% SP 800-115 (évaluation technique de la sécurité) ; guide du CHAI.
% Corrigés : corriges/Corr_C21.tex.
% ============================================================================

\chapter{Auditer la Sécurité Technique de l'Infrastructure}
\label{chap:audit-technique}
\chaptermark{Sécurité technique}

\epigraph{%
    « Une politique de sécurité dit ce qui devrait être fermé.\\
    Une configuration dit ce qui l'est. »%
}{Adage des auditeurs techniques}

\niveauChapitre{Intermédiaire à expert}{%
    Maîtrise du Chapitre~\ref{chap:audit-smsi} (SMSI, mesures) et des
    notions de réseau et de système du cursus d'ingénieur%
}

\begin{boxobjectifs}
\begin{itemize}
    \item Décomposer l'audit d'une infrastructure~: segmentation,
          durcissement, protocoles, vulnérabilités et correctifs,
          journalisation.
    \item Comparer les flux autorisés et les flux observés à partir d'un
          export de pare-feu.
    \item Vérifier un durcissement à partir d'un guide de configuration
          (CIS Benchmarks) et de commandes système.
    \item Auditer la gestion des vulnérabilités et des correctifs, et lire
          un score CVSS.
    \item Vérifier la journalisation centralisée.
    \item Coter par site avant d'agréger.
\end{itemize}
\end{boxobjectifs}

\bigskip

% ============================================================================
\section*{Question directrice : ce qui est écrit est-il configuré~?}
% ============================================================================

LiZbiZ exploite quatre centres de données, à Douala, Yaoundé, Garoua et
Bafoussam. Sa politique de sécurité interdit les protocoles
d'administration en clair, impose le changement des mots de passe par
défaut et sépare les postes des utilisateurs des serveurs de base de
données. Le domaine~D a montré que les accès étaient mal tenus
(Chapitre~\ref{chap:decomposition}). AUDIX CONSEIL doit maintenant dire si
l'infrastructure est configurée comme la politique l'écrit.

\medskip
\textbf{Plan du chapitre.}
\begin{enumerate}
    \item Décomposer l'audit de l'infrastructure (section~\ref{sec:c21-decomposer}).
    \item La segmentation~: flux autorisés, flux observés
          (section~\ref{sec:c21-flux}).
    \item Le durcissement et les protocoles (section~\ref{sec:c21-durcissement}).
    \item Les vulnérabilités et les correctifs (section~\ref{sec:c21-vuln}).
    \item La journalisation (section~\ref{sec:c21-journaux}).
    \item Quatre centres, une même politique (section~\ref{sec:c21-lizbiz}).
\end{enumerate}
Les tests actifs (scans, tests d'intrusion) et leur encadrement juridique
font l'objet du Chapitre~\ref{chap:intrusion}~; ce chapitre s'appuie sur
des exports et des lectures de configuration, sans action sur les systèmes.

\separateur

% ============================================================================
\section{Décomposer l'Audit de l'Infrastructure}
\label{sec:c21-decomposer}
% ============================================================================

\begin{tableaucours}{cL{3.0cm}L{5.6cm}L{2.2cm}}{%
    \tableauHeader{Code} &
    \tableauHeader{Sous-domaine} &
    \tableauHeader{Contrôle clé} &
    \tableauHeader{ISO/IEC 27002}
}
IN.1 & Architecture et segmentation & Seuls les flux autorisés par la
matrice de flux traversent les zones & 8.20, 8.22 \\
IN.2 & Durcissement & Aucun compte par défaut actif avec son mot de passe
d'origine & 8.9 \\
IN.3 & Protocoles & Aucun protocole d'administration ou de transfert en
clair & 8.20, 8.24 \\
IN.4 & Vulnérabilités et correctifs & Les correctifs critiques sont
appliqués dans le délai de la politique & 8.8 \\
IN.5 & Journalisation et surveillance & Les équipements envoient leurs
journaux à un point central surveillé & 8.15 \\
\end{tableaucours}
\vspace{-0.8em}
\begin{center}{\small\itshape Tableau~21.1. Décomposition de l'audit de
l'infrastructure (numéros de l'édition 2022, à vérifier)}\end{center}

IN.5 rejoint l'article~29 de la loi n°~2010/012 (surveillance des accès,
conservation des journaux, Chapitre~\ref{chap:droit}) et le sous-domaine
D.5 du contrôle d'accès~: l'auditeur réutilise les mêmes preuves.

\begin{boxinfo}
\textbf{Une lecture, pas une attaque.} Toutes les vérifications de ce
chapitre se font sur des exports produits par l'exploitant (règles de
pare-feu, fichiers de configuration, listes de comptes, état des
correctifs) ou par des commandes de lecture exécutées par un
administrateur, en présence de l'auditeur. Elles ne modifient rien et
n'exposent pas l'auditeur aux articles~65 et~69 de la loi n°~2010/012,
pourvu que la lettre de mission les prévoie.
\end{boxinfo}

% ============================================================================
\section{La Segmentation : Flux Autorisés, Flux Observés}
\label{sec:c21-flux}
% ============================================================================

\subsection{Deux matrices}

La segmentation se vérifie en comparant deux matrices~: celle des flux que
la politique autorise entre zones, et celle des flux que les règles des
pare-feu laissent réellement passer. L'écart entre les deux est le
constat.

\begin{tableaucours}{L{3.0cm}cccc}{%
    \tableauHeader{De \textbackslash{} vers} &
    \tableauHeader{DMZ} & \tableauHeader{Applicatif} &
    \tableauHeader{Utilisateurs} & \tableauHeader{Administration}
}
Internet & 443 & (aucun) & (aucun) & (aucun) \\
DMZ & (aucun) & 8443 & \textbf{445 (observé)} & (aucun) \\
Utilisateurs & 443 & 8443 \textbf{+ 1433 (observé)} & (aucun) & (aucun) \\
Administration & 22 & 22 & 22 & (aucun) \\
\end{tableaucours}
\vspace{-0.8em}
\begin{center}{\small\itshape Tableau~21.2. Matrice de flux d'un centre de
données~: ports autorisés, et en gras les flux observés non
autorisés}\end{center}

Dans cet exemple, deux flux non autorisés sont ouverts~: de la DMZ vers les
postes (un serveur web compromis pourrait atteindre les utilisateurs par
le partage de fichiers), et des postes vers la base de données (port~1433~:
un utilisateur peut contourner l'application et attaquer la base).

\subsection{Lire un export de règles}

\begin{boxoutils}{Repérer les flux non autorisés dans un export}
\begin{lstlisting}[language=Python]
import csv
# export du pare-feu : source, destination, port, action
# matrice autorisée : couples (zone source, zone destination, port)
autorises = {("internet", "dmz", 443), ("dmz", "app", 8443),
             ("users", "app", 8443), ("users", "dmz", 443),
             ("admin", "dmz", 22), ("admin", "app", 22),
             ("admin", "users", 22)}
with open("regles_garoua.csv", encoding="utf-8") as f:
    for r in csv.DictReader(f, delimiter=";"):
        cle = (r["zone_src"], r["zone_dst"], int(r["port"]))
        if r["action"] == "allow" and cle not in autorises:
            print("NON AUTORISE :", cle, "regle", r["id"])
\end{lstlisting}
Le test ne retient que les règles qui \emph{autorisent} (\texttt{allow}).
Une règle qui \emph{refuse} un flux non prévu n'est pas un écart~: c'est le
comportement attendu.
\end{boxoutils}

\begin{boxalert}
\textbf{Violation typique.} Compter comme écart une règle de refus, ou une
règle placée après une règle « tout refuser » qui ne sera jamais atteinte.

\smallskip
\textbf{Règle~:} un pare-feu applique ses règles dans l'ordre~; l'auditeur
reconstitue l'effet réel de la liste, et non la présence de chaque ligne.
\end{boxalert}

% ============================================================================
\section{Le Durcissement et les Protocoles}
\label{sec:c21-durcissement}
% ============================================================================

\subsection{Les guides de configuration}

\begin{boxdef}
\termecle{Durcissement}\index{durcissement}~: réduction de la surface
d'attaque d'un système par sa configuration~: services inutiles arrêtés,
comptes par défaut supprimés ou modifiés, paramètres de sécurité activés.
\end{boxdef}

Le Center for Internet Security publie des guides de configuration par
produit (systèmes, bases de données, navigateurs, équipements réseau),
les \emph{CIS Benchmarks}, et un ensemble de mesures prioritaires, les
\emph{CIS Controls}~\cite{cis2021controls}. L'organisation choisit sa
configuration de référence~; l'auditeur vérifie l'écart entre la référence
et l'état réel.

\subsection{Quelques vérifications en ligne de commande}

\begin{boxoutils}{Lecture d'un serveur Linux (exécutée par l'administrateur)}
\begin{lstlisting}[language=bash_forensique]
ss -tlnp                           # services en écoute et ports
awk -F: '$3 == 0 {print $1}' /etc/passwd   # comptes de super-utilisateur
grep -Ei '^(PermitRootLogin|PasswordAuthentication)' /etc/ssh/sshd_config
systemctl list-unit-files --state=enabled  # services activés au démarrage
\end{lstlisting}
Attendus de la politique de LiZbiZ~: aucun service Telnet ni FTP~; un seul
compte d'UID~0~; connexion directe de \texttt{root} interdite en SSH.
\end{boxoutils}

\begin{boxoutils}{Lecture d'un serveur Windows (PowerShell, administrateur)}
\begin{lstlisting}[language=bash_forensique]
Get-LocalUser | Select Name, Enabled, PasswordLastSet
Get-NetTCPConnection -State Listen | Select LocalPort, OwningProcess
auditpol /get /category:*          # stratégie d'audit des événements
\end{lstlisting}
Attendus~: compte « Administrateur » renommé ou désactivé~; aucun compte
« Invité » actif~; journalisation des ouvertures de session activée.
\end{boxoutils}

\begin{boxalert}
\textbf{Violation typique.} Exécuter soi-même des commandes sur un serveur
de production avec un compte prêté par l'administrateur.

\smallskip
\textbf{Règle~:} l'administrateur exécute, l'auditeur observe et reçoit la
sortie, horodatée et accompagnée de son empreinte. L'auditeur n'agit jamais
sur un système audité avec l'identité d'un autre~: ce serait une
usurpation, et la preuve serait contestable.
\end{boxalert}

\subsection{Les protocoles en clair}

Telnet, FTP, HTTP pour l'administration, SNMP en version~1 ou~2c
transmettent les identifiants en clair sur le réseau. Leur présence se lit
dans la liste des services en écoute et dans la configuration des
équipements. Le critère est simple~: aucun protocole d'administration en
clair~; une redirection de HTTP vers HTTPS n'est pas un écart si aucune
donnée sensible ne transite avant la redirection.

% ============================================================================
\section{Vulnérabilités et Correctifs}
\label{sec:c21-vuln}
% ============================================================================

\subsection{Le processus attendu}

La gestion des vulnérabilités suit un cycle~: veille (bulletins des
éditeurs), inventaire des systèmes concernés, évaluation de la gravité,
correction dans un délai fixé selon la gravité, vérification. L'affaire
Equifax (Chapitre~\ref{chap:defaillances}) montre ce que coûte une étape
manquante.

\subsection{Lire un score CVSS}

\begin{boxdef}
\termecle{CVSS}\index{CVSS} (\textit{Common Vulnerability Scoring
System})~: système de notation de la gravité des vulnérabilités, de 0 à 10,
publié par le FIRST~\cite{first2019cvss}. La note de base combine le vecteur
d'attaque (réseau, adjacent, local, physique), la complexité, les
privilèges requis, l'interaction de l'utilisateur, et l'effet sur la
confidentialité, l'intégrité et la disponibilité.
\end{boxdef}

\begin{tableaucours}{L{3.0cm}cL{6.0cm}}{%
    \tableauHeader{Gravité} &
    \tableauHeader{Note de base} &
    \tableauHeader{Délai de correction (politique de LiZbiZ)}
}
Critique & 9,0 à 10,0 & 30 jours (7 jours si exploitée activement) \\
Élevée & 7,0 à 8,9 & 30 jours \\
Moyenne & 4,0 à 6,9 & 90 jours \\
Faible & 0,1 à 3,9 & À la prochaine maintenance \\
\end{tableaucours}
\vspace{-0.8em}
\begin{center}{\small\itshape Tableau~21.3. Gravités CVSS et délais de
correction (les délais sont ceux de la politique de LiZbiZ)}\end{center}

\subsection{Tester le délai}

\begin{boxprocedure}
\textbf{Auditer le délai de correction}
\begin{enumerate}
    \item Obtenir la liste des correctifs critiques publiés pendant la
          période pour les systèmes du périmètre.
    \item Obtenir, pour chaque serveur, la date d'installation de chaque
          correctif (outil de gestion des correctifs, ou commande de
          lecture).
    \item Calculer, pour chaque couple serveur et correctif, le délai entre
          la publication et l'installation.
    \item Coter \NON{} au-delà du délai, \NEK{} si la date d'installation ne
          peut pas être établie.
\end{enumerate}
\end{boxprocedure}

\begin{boxalert}
\textbf{Violation typique.} Compter « en retard » un correctif installé le
trentième jour exactement.

\smallskip
\textbf{Règle~:} lire la politique~: « sous 30~jours » inclut le
trentième~; « en moins de 30~jours » ne l'inclut pas. L'auditeur écrit la
convention dans la fiche d'assertion avant le test.
\end{boxalert}

% ============================================================================
\section{La Journalisation}
\label{sec:c21-journaux}
% ============================================================================

Un journal qui reste sur la machine qu'il surveille peut être effacé par
celui qui l'a compromise. Le contrôle attendu est l'envoi en temps réel
vers un point central (serveur de journaux, outil de gestion des
événements de sécurité), protégé et surveillé.

\begin{tableaucours}{L{4.0cm}L{4.2cm}L{4.0cm}}{%
    \tableauHeader{Question} &
    \tableauHeader{Preuve} &
    \tableauHeader{Cote si absente}
}
Les équipements envoient-ils leurs journaux~? & Configuration de
l'envoi~; présence des journaux au point central & \NON{} \\
Les journaux sont-ils protégés~? & Droits sur le point central, empreintes
& \NON{} \\
Sont-ils surveillés~? & Règles d'alerte, traces de traitement des alertes &
\NON{} si aucune règle~; \NEK{} si non montré \\
Sont-ils conservés assez longtemps~? & Plus ancien journal disponible,
politique de conservation & \NON{} \\
\end{tableaucours}
\vspace{-0.8em}
\begin{center}{\small\itshape Tableau~21.4. Auditer la
journalisation}\end{center}

% ============================================================================
\section{Quatre Centres, une Même Politique}
\label{sec:c21-lizbiz}
% ============================================================================

\begin{boxlecture}{Les quatre centres de LiZbiZ}
AUDIX CONSEIL cote, pour chaque centre, le contrôle clé de chacun des cinq
sous-domaines, d'après les exports de configuration remis par la DSI.

\medskip
\begin{tableaucours}{L{2.6cm}cccccc}{%
    \tableauHeader{Centre} &
    \tableauHeader{IN.1} & \tableauHeader{IN.2} & \tableauHeader{IN.3} &
    \tableauHeader{IN.4} & \tableauHeader{IN.5} & \tableauHeader{O/N/K}
}
Douala & \NON{} & \OUI{} & \OUI{} & \OUI{} & \OUI{} & 4 / 1 / 0 \\
Yaoundé & \OUI{} & \NON{} & \OUI{} & \OUI{} & \OUI{} & 4 / 1 / 0 \\
Garoua & \NON{} & \NON{} & \NON{} & \NEK{} & \NON{} & 0 / 4 / 1 \\
Bafoussam & \NON{} & \NON{} & \NON{} & \NEK{} & \NEK{} & 0 / 3 / 2 \\
Ensemble & & & & & & 8 / 9 / 3 \\
\end{tableaucours}

\medskip
\textbf{Constatation.} Ensemble~: $\mu = 0{,}40$~; $U = 0{,}15$~;
$H \approx 1{,}46$~bit. Douala et Yaoundé ont chacun un écart~; Garoua et
Bafoussam n'appliquent pratiquement pas la politique. Pour ces deux
centres, l'état des correctifs n'a pas pu être établi~: aucun export n'a
été produit.

\medskip
\textbf{Analyse} (confiance moyenne). Non-conformité claire pour
l'ensemble, mais la cause est localisée~: selon la DSI, les centres de
Garoua et de Bafoussam ont été installés par un prestataire, sans
reprendre la configuration de référence (à confirmer par les dossiers
d'installation). Une moyenne des quatre centres serait trompeuse~; la
recommandation porte sur une configuration de référence imposée à chaque
nouveau site, et vérifiée avant sa mise en service.

\medskip
\hfill\lecture{0{,}40}{0{,}15}{1{,}46}
\end{boxlecture}

\begin{boxinfo}
\textbf{Décomposer l'entropie.} Avec $U = 0{,}15$, la part d'ignorance
vaut $H_b(0{,}15) \approx 0{,}61$~bit, soit environ 42~\% de $H$~; le reste
vient du désaccord entre \OUI{} et \NON{}, c'est-à-dire de la différence
entre les deux centres conformes et les deux autres
(Chapitre~\ref{chap:conclure}). La décomposition par site confirme ce que
dit l'entropie.
\end{boxinfo}

\begin{boxmetier}{Auditeur de sécurité}
\textbf{Lire un export de configuration}
\begin{enumerate}
    \item Vérifier la date, l'équipement et l'empreinte de l'export.
    \item Lister les services exposés et les comparer à la configuration de
          référence.
    \item Chercher les comptes par défaut et les protocoles en clair
          (Telnet, FTP, HTTP d'administration).
    \item Reconstituer les flux autorisés par les règles, dans leur ordre,
          et les comparer à la matrice de flux.
    \item Coter chaque contrôle par centre, avant toute agrégation.
\end{enumerate}
\end{boxmetier}

\begin{boxmetier}{RSSI}
\textbf{Préparer un audit technique}
\begin{enumerate}
    \item Fournir la matrice de flux autorisés et la configuration de
          référence de chaque type d'équipement.
    \item Produire les exports demandés, avec leur date et leur empreinte.
    \item Désigner un administrateur par site pour exécuter les commandes
          de lecture devant l'auditeur.
    \item Traiter les constats par ordre de gravité, et documenter chaque
          correction pour l'audit de suivi.
\end{enumerate}
\end{boxmetier}

\separateur

\begin{boxresume}
\textbf{Ce qu'il faut retenir de ce chapitre~:}
\begin{itemize}
    \item \textbf{Cinq sous-domaines d'infrastructure}~: segmentation,
          durcissement, protocoles, vulnérabilités, journalisation.
    \item \textbf{La segmentation se prouve} en comparant flux autorisés et
          flux observés, en lisant les règles dans leur ordre.
    \item \textbf{Le durcissement se mesure} contre une configuration de
          référence (CIS Benchmarks)~; l'administrateur exécute, l'auditeur
          observe.
    \item \textbf{Les correctifs} se jugent sur un délai lié à la gravité
          CVSS, écrit dans la politique~; une date inconnue est \NEK{}.
    \item \textbf{Les journaux} valent s'ils sont centralisés, protégés,
          surveillés et conservés.
    \item \textbf{Coter par site avant d'agréger}~: une moyenne peut cacher
          deux centres hors politique.
\end{itemize}
\end{boxresume}

\begin{boxinfo}
\textbf{Points de contrôle de l'élève.} Je sais~: dresser une matrice de
flux autorisés et observés~; lire un export de règles dans l'ordre~;
interpréter les sorties des commandes de lecture~; associer une note CVSS
à une gravité et un délai~; auditer la journalisation~; coter par site.
\end{boxinfo}

\bigskip

\begin{boxexo}[title={\ding{45}\quad Exercice~21.1. Rattacher \niveaubase}]
Rattachez à IN.1 à IN.5~: (a) SNMP v2c activé avec la communauté
« public »~; (b) aucun envoi des journaux du pare-feu de Garoua~; (c) le
compte « admin » du routeur a son mot de passe d'usine~; (d) les postes
peuvent joindre la base en 1433~; (e) un correctif critique installé après
75~jours.
\end{boxexo}

\begin{boxexo}[title={\ding{45}\quad Exercice~21.2. Matrice de flux \niveaumoyen}]
La politique de LiZbiZ autorise~: Internet vers DMZ en 443~; DMZ vers
applicatif en 8443~; utilisateurs vers applicatif en 8443~; administration
vers toutes les zones en 22. L'export du pare-feu de Bafoussam montre en
outre~: utilisateurs vers applicatif en 1433 et en 3389~; Internet vers
administration en 23.
\begin{enumerate}[(a)]
    \item Dressez la matrice des flux autorisés et observés.
    \item Classez les écarts par gravité, en justifiant.
    \item Quelle cote retenir pour IN.1 et pour IN.3 à Bafoussam~?
    \item Rédigez la constatation de l'observation la plus grave.
\end{enumerate}
\end{boxexo}

\begin{boxexo}[title={\ding{45}\quad Exercice~21.3. L'ordre des règles \niveaumoyen}]
Un pare-feu applique~: (1) refuser tout de « users » vers « app » en 1433~;
(2) autoriser tout de « users » vers « app »~; (3) refuser tout. Le flux
1433 des utilisateurs vers la base passe-t-il~? Et le flux 3389~? Que
devient le résultat si l'on inverse les règles 1 et 2~?
\end{boxexo}

\begin{boxexo}[title={\ding{45}\quad Exercice~21.4. Interpréter des sorties \niveaumoyen}]
Sur un serveur Linux de Garoua, \texttt{ss -tlnp} montre les ports 21, 22,
23 et 3306 en écoute~; \texttt{awk} sur \texttt{/etc/passwd} rend
\texttt{root} et \texttt{support}. Quels constats faites-vous, quelles cotes
pour IN.2 et IN.3, et quelles vérifications complémentaires demandez-vous~?
\end{boxexo}

\begin{boxexo}[title={\ding{45}\quad Exercice~21.5. CVSS et délais \niveaubase}]
Trois vulnérabilités sont publiées le 1\textsuperscript{er}~mars~2026~:
note 9,8 (exploitée activement), 7,5 et 5,3. Les correctifs sont installés
à Douala le 6~mars, le 31~mars et le 15~mai. Cotez chaque correctif selon le
Tableau~21.3.
\end{boxexo}

\begin{boxexo}[title={\ding{45}\quad Exercice~21.6. Le correctif du trentième jour \niveaumoyen}]
La politique dit « dans un délai de 30~jours ». Un correctif publié le
2~avril est installé le 2~mai. Est-il dans le délai~? Comment
l'auditeur évite-t-il que la réponse dépende de son humeur~?
\end{boxexo}

\begin{boxexo}[title={\ding{45}\quad Exercice~21.7. La journalisation \niveaumoyen}]
À Yaoundé, les journaux sont envoyés au serveur central, mais
l'administrateur système y a les droits d'écriture et de suppression.
Cotez les quatre questions du Tableau~21.4 et rédigez la recommandation.
\end{boxexo}

\begin{boxexo}[title={\ding{45}\quad Exercice~21.8. Le script de flux \niveaumoyen}]
Modifiez le script de la section~\ref{sec:c21-flux} pour qu'il tienne compte
de l'ordre des règles (la première règle qui correspond s'applique).
Décrivez votre algorithme, puis écrivez-le.
\end{boxexo}

\begin{boxexo}[title={\ding{45}\quad Exercice~21.9. Une configuration de référence \niveauexpert}]
Rédigez les dix contrôles de la configuration de référence qu'AUDIX
CONSEIL recommande d'imposer à tout nouveau centre de LiZbiZ, chacun avec
son critère et la commande ou l'export qui permet de le vérifier.
\end{boxexo}

\begin{boxexo}[title={\ding{45}\quad Exercice~21.10. Agréger ou non \niveauexpert}]
Le directeur général veut une seule ligne pour l'infrastructure dans la
synthèse. Rédigez-la à partir de la lecture des quatre centres, sans
cacher la différence entre les sites, et proposez la priorité des
recommandations.
\end{boxexo}

\bigskip

\begin{boxinfo}
\textbf{Pour aller plus loin.}
\begin{itemize}
    \item Center for Internet Security, \textit{CIS Controls}
          v8~\cite{cis2021controls} et les CIS Benchmarks.
    \item FIRST, \textit{CVSS v3.1 Specification}~\cite{first2019cvss}.
    \item NIST, SP~800-115~\cite{nist2008sp800115}~: le guide de
          l'évaluation technique de la sécurité.
    \item ISO/IEC~27002:2022~\cite{iso27002_2022}~: les mesures
          technologiques.
\end{itemize}
\end{boxinfo}

\begin{boxinfo}
\textbf{Transition.} La lecture des configurations montre ce qui est
paramétré~; elle ne montre pas ce qu'un attaquant pourrait en faire. Le
Chapitre~\ref{chap:intrusion} présente les tests d'intrusion et l'audit
de code, leur cadre légal et leur rapport.
\end{boxinfo}
